\section{Compact Prompt 与 Handoff Prompt 的内容}

\subsection{压缩提示词（Compact Prompt）}

Codex 源码中为非 OpenAI Provider 提供了明文的压缩提示词（\texttt{compact.prompt.md}），核心指令如下：

\begin{importantbox}{Compact Prompt 全文要义}
"你正在执行一个上下文检查点的压缩。请为另一个将接手该任务的 LM 创建一份交接摘要，包含：
\begin{itemize}
    \item 当前的进度
    \item 已经做出的关键决策
    \item 重要的上下文、约束条件和用户偏好
    \item 尚未完成的工作
    \item 继续工作所需的任何关键数据、示例或资料
\end{itemize}
一定要简洁、结构化，并专注于帮助下一个 LM 无缝衔接工作。"
\end{importantbox}

\subsection{交接提示词（Handoff Prompt / Summary Prefix）}

交接发生在压缩完成之后，\texttt{summary.prefix.md} 的内容如下：

\begin{importantbox}{Handoff Prompt 全文要义}
"另外一个语言模型已经开始解决这个问题，并生成了其思考过程的摘要。你也可以访问该语言模型所使用的工具状态。请利用这些信息在已有工作的基础上继续进行，避免重复劳动。以下就是该语言模型的摘要——请利用摘要中的信息辅助你自己的分析。"
\end{importantbox}

\subsection{完整的两步流程}

\begin{enumerate}
    \item \textbf{Step 1 — 压缩}：
    \begin{itemize}
        \item 输入：System Prompt + Compact Prompt + 全部对话历史。
        \item 输出：一份结构化的"密文"（压缩摘要）。
    \end{itemize}
    \item \textbf{Step 2 — 交接}：
    \begin{itemize}
        \item 输入：System Prompt + Handoff Prompt + 压缩摘要 + 用户的新请求。
        \item 输出：正常的 Agent 响应，无缝继续工作。
    \end{itemize}
\end{enumerate}

\subsection{本章小结}

Compact Prompt 指导模型生成"交接摘要"，Handoff Prompt 指导接手模型利用摘要继续工作。两者配合实现了"压缩$\rightarrow$交接$\rightarrow$继续"的闭环。


